\begin{frame}{자주 묻는 질문 (5.3) --- 단일 샘플 vs 전체 궤적}
  \textbf{``단일 샘플''과 ``전체 궤적'' 중요도 샘플링의
  차이는?}
  \begin{itemize}
    \item 위에서 쓰던 $\rho = \prod_t \frac{\pi(a_t|s_t)}{b(a_t|s_t)}$
      은 \kb{전체 궤적 (full trajectory)} 중요도 --- 에피소드
      전체 확률의 비율.
    \item 상태가치 $V^\pi(s)$ 를 추정할 때는 \kb{상태 $s$ 이후
      부분}(partial trajectory)만 쓰면 충분 --- 분모가 ``$s$
      에서 시작하는 궤적 확률의 비율''이라 전체 궤적보다
      \textbf{분산이 작다}.
    \item $Q^\pi(s,a)$ 를 추정할 때는 해당 $(s,a)$ 이후 부분만
      쓴다.
  \end{itemize}
  \vskip0.5em
  \begin{block}{실전 원칙}
    \textbf{가능한 한 짧은(목표와 관련된) 부분 궤적만 가중치로
    쓴다} --- 이것이 분산을 줄이는 가장 간단한 방법.
    (5.2절 블랙잭에서 ``플레이어 합 12 이상 상태에서만 $Q$ 를
    학습''한 것도 같은 원리 --- 시작 상태부터가 아니라
    \textbf{관심 상태 이후}의 리턴만 쓴다.)
  \end{block}
\end{frame}
